\chapter{Dry Skin (Xerosis)}

\section*{Definition}
A common skin condition characterized by rough, scaling, flaking, and sometimes pruritic skin due to reduced water content in the stratum corneum, often exacerbated by cold weather, low humidity, and frequent washing.

\section*{Pathophysiology}
Dry skin results from disruption of the skin barrier function, particularly the stratum corneum lipid matrix (ceramides, cholesterol, free fatty acids) that normally retains water. Decreased natural moisturizing factors (amino acids, urea, lactate) and impaired desquamation lead to accumulation of corneocytes on the surface. Environmental factors (low humidity, cold, wind, excessive bathing) accelerate transepidermal water loss. Aging reduces sebaceous and sweat gland activity.

\section*{Causes}
\begin{itemize}
  \item Cold, dry weather (winter xerosis, most common)
  \item Excessive bathing or hand washing (strips natural oils)
  \item Harsh soaps and detergents (disrupt skin barrier)
  \item Aging (decreased sebum production, thinner skin)
  \item Central heating or air conditioning (low indoor humidity)
  \item Atopic dermatitis (eczema, genetic barrier defect)
  \item Psoriasis (increased epidermal turnover with scaling)
  \item Hypothyroidism (reduced sweating and sebum production)
  \item Diabetes mellitus (dehydration, autonomic dysfunction)
  \item Medications (diuretics, retinoids, antihistamines, statins)
\end{itemize}

\section*{When to See a Doctor}
See your doctor if:
\begin{itemize}
  \item Severe or worsening symptoms
  \item Dry skin with severe itching, signs of infection (increasing redness, warmth, swelling, pus), extensive cracking or fissures, or rash unresponsive to moisturizers
  \item Accompanied by other concerning signs
\end{itemize}

\section*{Related Symptoms}
\begin{itemize}
  \item Itching (pruritus)
  \item Rash
  \item Skin flaking
  \item Erythema
  \item Skin fissures
\end{itemize}
\textbf{Important:} If this symptom persists, your doctor will check for serious underlying causes.

\section*{Self-Care / Home Management}
\begin{itemize}
  \item Rest and avoid activities that make it worse.
  \item Maintain adequate hydration and a balanced diet.
  \item Over-the-counter remedies may provide symptomatic relief (consult a pharmacist or doctor).
  \item Monitor symptoms and seek professional advice if they persist or worsen.
  \item Avoid self-medicating with prescription medications without medical guidance.
\end{itemize}

\section*{How Your Doctor Will Evaluate You}
\begin{itemize}
  \item A thorough history and physical examination are the first steps in evaluation.
  \item Laboratory tests (blood work, urinalysis) may be ordered based on clinical suspicion.
  \item Imaging studies (X-ray, ultrasound, CT, MRI) may be indicated when structural pathology is suspected.
  \item Specialist referral is arranged when the diagnosis remains uncertain or requires specific expertise.
\end{itemize}

\section*{Treatment Options}
\begin{itemize}
  \item Treatment targets the underlying cause identified through diagnostic evaluation.
  \item Supportive care and symptom management are provided while awaiting definitive therapy.
  \item Medications, lifestyle modifications, and physical therapies may be prescribed as appropriate.
  \item Follow-up ensures treatment effectiveness and allows timely adjustments to the care plan.
\end{itemize}

\clearpage